% =====================================================================
%  Dokumen LaTeX lengkap (standalone) — siap copy-paste ke Overleaf.
%  Berisi tabel perbandingan fitur: penelitian terkait vs penelitian ini.
%
%  Cara pakai di Overleaf:
%   1. Buat project baru, pilih compiler pdfLaTeX.
%   2. Copy-paste seluruh isi berkas ini ke main.tex, lalu Recompile.
%
%  Untuk menyisipkan tabelnya saja ke naskah IJECE Anda, salin hanya
%  blok di antara "=== TABEL MULAI ===" dan "=== TABEL SELESAI ===".
%  Pastikan preamble naskah memuat: booktabs, tabularx, amssymb.
% =====================================================================
\documentclass[10pt,a4paper]{article}

\usepackage[utf8]{inputenc}
\usepackage[T1]{fontenc}
\usepackage[margin=2cm]{geometry}
\usepackage{booktabs}     % \toprule \midrule \bottomrule
\usepackage{tabularx}     % kolom lebar otomatis
\usepackage{amssymb}      % \checkmark
\usepackage{array}

\begin{document}

% =====================================================================
%  === TABEL MULAI ===
%  (Dalam naskah dua kolom seperti IJECE, ganti {table} -> {table*}
%   agar tabel melebar penuh melintasi dua kolom.)
% =====================================================================
\begin{table}[htbp]
\caption{Perbandingan fitur antara penelitian terkait dan penelitian yang diusulkan}
\label{tab:feature_comparison}
\centering
\footnotesize
\setlength{\tabcolsep}{4pt}
\renewcommand{\arraystretch}{1.3}
\begin{tabularx}{\textwidth}{@{} p{2.9cm} c c c c c c c @{}}
\toprule
\textbf{Studi} &
\textbf{\shortstack{Air-gapped\\backup}} &
\textbf{\shortstack{Kompresi\\multi-algoritma}} &
\textbf{\shortstack{Seleksi\\adaptif}} &
\textbf{\shortstack{Enkripsi\\AES-256}} &
\textbf{\shortstack{Deteksi\\anomali}} &
\textbf{\shortstack{Recovery\\otomatis}} &
\textbf{\shortstack{Integrasi\\end-to-end}} \\
\midrule
Smith \& Lee (2023)          & --          & \checkmark & --          & --          & --          & --          & -- \\
Kumar \& Patel (2022)        & --          & \checkmark & --          & --          & --          & --          & -- \\
Novak et al. (2021)          & \checkmark  & --          & --          & --          & \checkmark & \checkmark  & -- \\
Sahu \& Panda (2025)         & --          & \checkmark & --          & --          & --          & --          & -- \\
Thomas et al. (backup)       & \checkmark  & --          & --          & --          & --          & \checkmark  & -- \\
Kim et al. (CryptoSniffer)   & --          & --          & --          & --          & \checkmark & --          & -- \\
Kumar et al. (XGBoost)       & --          & --          & --          & --          & \checkmark & --          & -- \\
NIST SP 800-209 (2024)       & \checkmark  & --          & --          & --          & --          & --          & -- \\
\midrule
\textbf{Penelitian ini}      & \checkmark  & \checkmark & \checkmark  & \checkmark  & \checkmark & \checkmark  & \checkmark \\
\bottomrule
\end{tabularx}

\vspace{4pt}
\raggedright\footnotesize
\textit{Keterangan}: \checkmark{} = fitur tersedia; -- = tidak tersedia/tidak dibahas.
Kolom \textit{Seleksi adaptif} (pemilihan algoritma kompresi berbasis aturan menurut tipe
dan ukuran berkas) serta \textit{Integrasi end-to-end} (menggabungkan air-gapped, kompresi
adaptif, enkripsi AES-256, deteksi anomali, dan recovery otomatis dalam satu kerangka kerja)
merupakan fitur yang \textbf{hanya dimiliki oleh penelitian ini} dan menjadi pembeda utama
(\textit{novelty}) dari studi terdahulu.
\end{table}
% =====================================================================
%  === TABEL SELESAI ===
% =====================================================================

\end{document}
